% ============================================================
% Выходные сведения (оборот титула)
% ============================================================
\pagestyle{fancy}
\fancyhf{}
\renewcommand{\headrulewidth}{1pt}
\renewcommand{\footrulewidth}{1pt}

\lhead{\color{black}\devicecode}
\chead{\color{black}Редакция \rev~от~\revdate}
\rhead{\color{unigreen}\devicetype}
\lfoot{\color{unigreen}Руководство по эксплуатации}
\rfoot{\thepage}

\pagecolor{white}

% Блок с копирайтом, редакцией и датой
\noindent
$\copyright$ \the\year{} Юнител Инжиниринг
\hfill
\parbox{0.3\textwidth}{\raggedleft Редакция: \textcolor{unigray}{\rev}}

\vspace{3mm}
\noindent
Москва
\hfill
\parbox{0.3\textwidth}{\raggedleft Дата: \textcolor{unigray}{\revdate}}
\vspace{8mm}

% Основной текст (абзацы разделены пустыми строками)
Настоящее руководство по эксплуатации (РЭ) распространяется на шкаф электротехнический автоматики управления выключателем 110–220 кВ, УРОВ типа «\devicetype» (именуемый в дальнейшем «шкаф») и содержит сведения о его технических характеристиках, правилах эксплуатации, назначении и обслуживанию.

Необходимые сведения по функциональному назначению, основным параметрам, принципам работы, характеристикам, функциональным схемам формирования сигналов, перечню уставок и настраиваемых параметров приведены в руководствах по эксплуатации на соответствующие микропроцессорные устройства серии «ЮНИТ-М500», входящие в состав шкафа.

Настоящее руководство разработано в соответствии с требованиями технических условий \ManualTechSpec.

До включения шкафа в работу необходимо ознакомиться с настоящим руководством и руководством по эксплуатации на микропроцессорное устройство, поставляемое в комплекте со шкафом, и с эксплуатационной документацией на отдельные элементы шкафа, входящие в его состав.

Надежность и долговечность шкафа обеспечиваются не только его качеством, но и соблюдением режимов, условий эксплуатации, требований по транспортированию, хранению, монтажу. Выполнение всех требований, изложенных в настоящем документе, является обязательным.

К эксплуатации шкафа допускаются лица, изучившие настоящее руководство, прошедшие проверку знаний правил техники безопасности и эксплуатации электроустановок, электрических станций и подстанций, и имеющие группу по электробезопасности не ниже III по работе с электроустановками напряжением до 1000~В.

\vspace{35mm}
Компания Юнител Инжиниринг оставляет за собой авторские права на данный документ и на информацию, содержащуюся в нем, включая права на использование патентов. Копирование, использование и передача информации третьим лицам без письменного разрешения компании категорически запрещены.

\vspace{10mm}
По всем интересующим вопросам можно обращаться:
\begin{itemize}
  \item по адресу электронной почты: \href{mailto:rza@uni-eng.ru}{\color{unigreen}\uline{rza@uni-eng.ru}};
  \item по телефону: +7 (495) 165-99-98 (доб. 2561).
\end{itemize}
Также с нашей продукцией можно ознакомиться на сайте: \href{https://uni-eng.ru}{\color{unigreen}\uline{https://uni-eng.ru}}.